\begin{figure}[ht]
\usetikzlibrary{decorations.markings, arrows.meta}
\definecolor{kkcontour}{RGB}{216,90,48}
\definecolor{kkpole}{RGB}{226,75,74}
\definecolor{kkshade}{RGB}{55,138,221}
  \centering
  \begin{tikzpicture}[>=Stealth, scale=0.95,
      contourarrow/.style={postaction={decorate},
        decoration={markings,
          mark=at position 0.13 with {\arrow{Stealth[length=2.4mm]}},
          mark=at position 0.30 with {\arrow{Stealth[length=2.0mm]}},
          mark=at position 0.36 with {\arrow{Stealth[length=2.4mm]}},
          mark=at position 0.70 with {\arrow{Stealth[length=2.4mm]}}}}]

    % regione di analiticità (semidisco superiore)
    \fill[kkshade!8] (4,0) arc (0:180:4) -- cycle;

    % assi
    \draw[->, gray!60, thin] (-4.9,0) -- (5.0,0) node[below right, black]{$\operatorname{Re}\omega$};
    \draw[->, gray!60, thin] (0,-0.7) -- (0,5.0) node[left, black]{$\operatorname{Im}\omega$};

    % contorno di integrazione (asse reale + indentazione + arco grande), orientato in senso antiorario
    \draw[very thick, kkcontour, contourarrow]
          (-4,0) -- (1.65,0)
          arc (180:0:0.35)      % indentazione sopra il polo, da sinistra a destra
          -- (4,0)
          arc (0:180:4);        % semicerchio grande, da destra a sinistra

    % polo
    \fill[kkpole] (2,0) circle (2.2pt);
    \node[below=2pt] at (2,0) {$\omega_0$};
    \node[below left=-1pt] at (0,0) {$O$};

    % etichette
    \node[kkcontour] at (3.25,3.05) {$\Gamma$};
    \node[kkcontour, right] at (2.25,0.62) {$C_\eps$};
    \node[align=center] at (-1.7,2.0)
          {$\chi(\omega)$ olomorfa};

  \end{tikzpicture}
  \caption{Dominio di integrazione di $\chi(\omega)/(\omega-\omega_0)$ per le relazioni di
  Kramers--Kronig. Il contorno percorre l'asse reale, aggira il polo semplice in $\omega_0$
  con la semicirconferenza $C_\eps$ e si chiude con il semicerchio $\Gamma$ nel
  semipiano superiore, dove $\chi$ è olomorfa per causalità. L'orientazione è antioraria.}
  \label{fig:contorno-kk}
\end{figure}